\documentclass[12pt]{article}
\usepackage[margin=0.6in,top=1.15in,bottom=0.55in]{geometry}
\usepackage{lmodern}
\usepackage{microtype}
\usepackage{titlesec}
\usepackage{enumitem}
\usepackage[hidelinks]{hyperref}
\usepackage{../../latex/grader-worksheet}

\setlength{\parindent}{0pt}
\setlength{\parskip}{0.25em}
\titleformat{\section}{\large\bfseries\scshape}{}{0pt}{}

\GraderSetup{assignment-id=U1L5LE,template-revision=1}

\begin{document}

\begin{center}
{\LARGE\bfseries Week 5 Lit Example Worksheet}\par\vspace{0.05em}
{\large\scshape Julius Caesar: Brutus, Antony, and Formal Shape}\par\vspace{0.12em}
\rule{0.78\textwidth}{0.5pt}
\end{center}

\textbf{Purpose:} Apply Week 5's Whately method to Julius Caesar. Identify mood and formal validity before judging factual strength.

\textbf{Directions:} Use the Lit Example Reader. Your answers should distinguish form test from evidence test.

\section*{Part 1: Brutus's Form}

\textbf{Question 1.1}\\[-0.15em]
Reconstruct Brutus's argument in three claims: major premise, minor premise, and conclusion.

\GraderAnswerBox{Part 1 Q1}{2.45in}

\textbf{Question 1.2}\\[-0.15em]
State the mood of your three claims (A, E, I, or O for each). Is the form valid if written as All M are P; All S are M; therefore All S are P?

\GraderAnswerBox{Part 1 Q2}{2.25in}

\newpage

\section*{Part 2: Middle Term and Figure}

\textbf{Question 2.1}\\[-0.15em]
Identify S, P, and M in Brutus's formal map. Then write the figure pattern using S, P, and M.

\GraderAnswerBox{Part 2 Q1}{2.35in}

\textbf{Question 2.2}\\[-0.15em]
Does Brutus's formal shape by itself prove that Caesar was dangerously ambitious? Explain the difference between valid form and proved premise.

\GraderAnswerBox{Part 2 Q2}{2.45in}

\newpage

\section*{Part 3: Antony's Reply}

\textbf{Question 3.1}\\[-0.15em]
List Antony's three pieces of evidence against the claim that Caesar was dangerously ambitious.

\GraderAnswerBox{Part 3 Q1}{2.25in}

\textbf{Question 3.2}\\[-0.15em]
Does Antony give a cleaner formal syllogism than Brutus, or does he mainly attack Brutus's minor premise? Explain.

\GraderAnswerBox{Part 3 Q2}{2.45in}

\newpage

\section*{Part 4: Formal Judgment}

\textbf{Question 4.1}\\[-0.15em]
Write one possible syllogism for Antony's reasoning. Then name its mood if possible and say whether the conclusion follows from your premises.

\GraderAnswerBox{Part 4 Q1}{2.75in}

\textbf{Question 4.2}\\[-0.15em]
Which speaker gives the stronger form? Which gives stronger evidence against the pressure-point premise? Explain the difference.

\GraderAnswerBox{Part 4 Q2}{2.25in}

\end{document}
